\input{../tex-inputs/latex-leseansicht-vorspann}

\section[Gustav Schwarzkopf an Arthur Schnitzler, 14. 4. 1894]{L04320 Gustav Schwarzkopf an Arthur Schnitzler, 14. 4. 1894}
\nopagebreak\mylabel{L04320v}
\rehead{ }\normalsize\beginnumbering\briefempfaengerindex{Schnitzler, Arthur@\textsc{Schnitzler, Arthur}!zzzSchwarzkopf, Gustav@\emph{von Gustav Schwarzkopf}!1894-04-142@{14. 4. 1894}|(be}
\toendnotes[C]{\smallbreak\pagebreak[2]}
\correspDesc{Versand  durch Gustav Schwarzkopf am 14. 4. 1894 in Wien
\newline{}Erhalt  durch Arthur Schnitzler am [14. 4. 1894?] in Wien}\toendnotes[C]{\smallbreak}
\Standort{CUL, Schnitzler, B 96a.}
\physDesc{Briefkarte, 239 Zeichen
\newline{}Handschrift: schwarze Tinte, deutsche Kurrent}
\pstart
           \raggedleft{}{\pb}14. IV 94\pend
           
\pstart{}Verehrteſter Freund!\pend\vspace{0.5em}
\pstart
           Ich habe den So{\geminationn}tag{ }ſchon beſetzt von drei
               Uhr Nachmittags vor. Mit einer Kegelpartie. Ka{\geminationn} alſo am
               Vormittag nichts Größeres unternehmen. Schade, daß Sie gerade \uline{diesen} So{\geminationn}tag wählen mußten.\pend
           
\pstart
           Herzlichſt Ihr{\\[\baselineskip]}\spacefill\mbox{GustavSchwzk}\pend
           \leftskip=0em{}\selectlanguage{ngerman}\endnumbering\briefempfaengerindex{Schnitzler, Arthur@\textsc{Schnitzler, Arthur}!zzzSchwarzkopf, Gustav@\emph{von Gustav Schwarzkopf}!1894-04-142@{14. 4. 1894}|)be}\mylabel{L04320h}
\begin{anhang}
\end{anhang}\newcommand{\dateiname}{L04320}\newcommand{\titel}{Gustav Schwarzkopf an Arthur Schnitzler, 14. 4. 1894}\newcommand{\editorInnen}{Herausgegeben von Jahnke, SelmaMüller, Martin Anton}\input{../tex-inputs/latex-leseansicht-abspann}
